\documentclass[aspectratio=169]{beamer}

\usetheme{Madrid}
\usecolortheme{default}
\setbeamertemplate{navigation symbols}{}

\usepackage[utf8]{inputenc}
\usepackage[T1]{fontenc}
\usepackage[french]{babel}
\usepackage{amsmath, amssymb}
\usepackage{booktabs}
\usepackage{mathtools}

\title[Finsler-MDS]{Algorithmes d'optimisation pour Finsler-MDS}
\subtitle{Métriques asymétriques et distances géodésiques dans l'espace d'embedding}
\author{}
\date{}

\newcommand{\R}{\mathbb{R}}
\newcommand{\Dtarget}{\Delta}
\newcommand{\Fmetric}{F_\theta}
\newcommand{\stress}{\mathcal{S}}

\begin{document}

\begin{frame}
  \titlepage
\end{frame}

\begin{frame}{Objectif commun}
  On cherche un embedding $X=(x_1,\ldots,x_n)\in \R^{n\times d}$ tel que
  les distances dans l'embedding reproduisent des dissimilarités dirigées
  $\Dtarget_{ij}$.

  \vspace{0.4cm}
  \[
    \stress(X)
    =
    \sum_{(i,j)\in \mathcal{P}}
    w_{ij}
    \left(d_X(i,j)-\Dtarget_{ij}\right)^2 .
  \]

  \vspace{0.2cm}
  Deux choix principaux:
  \begin{itemize}
    \item distance directe: $d_X(i,j)=\Fmetric(x_j-x_i)$;
    \item distance géodésique: $d_X(i,j)$ est mesurée le long du graphe kNN de l'embedding.
  \end{itemize}
\end{frame}

\begin{frame}{Vue d'ensemble des optimiseurs}
  \begin{table}
    \centering
    \begin{tabular}{lll}
      \toprule
      Optimiseur & Distance optimisée & Idée principale \\
      \midrule
      Randers-SMACOF & directe & majoration quadratique dédiée à Randers \\
      Gradient descent & directe & baseline générique, gradients explicites \\
      Path-frozen & géodésique dure & Dijkstra puis chemins gelés \\
      Soft-Bellman-Ford & géodésique soft & relaxations différentiables par softmin \\
      DataSP / soft-FW & all-pairs soft & soft Floyd-Warshall différentiable \\
      \bottomrule
    \end{tabular}
  \end{table}

  \vspace{0.2cm}
  En pratique, le coût vient moins de la formule du stress que du choix de
  $\mathcal{P}$, du graphe de support, et de la façon dont les chemins sont
  recalculés.
\end{frame}

\begin{frame}{Randers-SMACOF}
  Principe: adapter SMACOF à une métrique de Randers
  \[
    F(u)=\|u\|+\alpha \langle v,u\rangle .
  \]

  \vspace{0.3cm}
  Améliorations implémentées:
  \begin{itemize}
    \item warm-start explicite: UMAP, Isomap, ancien embedding, etc.;
    \item support des poids $w_{ij}$ et du stress normalisé;
    \item solveurs linéaires adaptés au cas dense;
    \item backend GPU pour le cas Randers dense, surtout utile quand $n$ est grand;
    \item parallélisation des initialisations multiples avec \texttt{n\_jobs}.
  \end{itemize}

  \vspace{0.2cm}
  Rôle typique: obtenir vite une géométrie globale raisonnable avant de passer
  aux distances géodésiques.
\end{frame}

\begin{frame}{Pourquoi passer aux distances géodésiques ?}
  Avec des distances directes, l'embedding est poussé à représenter les
  dissimilarités par des segments droits.

  \vspace{0.3cm}
  Pour des données sur une variété, on préfère:
  \[
    d_X(i,j)
    \approx
    \min_{\gamma:i\to j}
    \sum_{(a,b)\in \gamma}
    \Fmetric(x_b-x_a).
  \]

  \vspace{0.2cm}
  Difficultés:
  \begin{itemize}
    \item le graphe kNN dépend de $X$;
    \item le plus court chemin change de manière discrète;
    \item all-pairs devient trop coûteux pour $n$ de quelques milliers;
    \item les gradients doivent remonter jusqu'aux coordonnées des arêtes.
  \end{itemize}
\end{frame}

\begin{frame}{Path-frozen}
  Principe:
  \begin{enumerate}
    \item reconstruire le graphe kNN depuis l'embedding courant;
    \item lancer Dijkstra depuis les sources actives;
    \item geler les chemins trouvés;
    \item optimiser plusieurs pas avec ces chemins fixes;
    \item recommencer.
  \end{enumerate}

  \vspace{0.2cm}
  Avec les chemins gelés,
  \[
    d_X(i,j)
    =
    \sum_{(a,b)\in \gamma_{ij}^{\mathrm{frozen}}}
    \Fmetric(x_b-x_a)
  \]
  devient différentiable et peu coûteux à réévaluer.
\end{frame}

\begin{frame}{Path-frozen: heuristiques importantes}
  \begin{itemize}
    \item \textbf{Outer/inner loop}: Dijkstra est relancé seulement aux outer iterations;
    \item \textbf{landmark\_mode="sources"}: peu de parcours de graphe, beaucoup de cibles;
    \item \textbf{max\_targets\_per\_source}: limite le nombre de paires optimisées par source;
    \item \textbf{target\_sampling}: random, farthest ou mixed selon le compromis global/local;
    \item \textbf{local\_neighbors}: ajoute des contraintes locales issues de $\Dtarget$;
    \item \textbf{local\_pair\_mode}: local direct ou géodésique.
  \end{itemize}

  \vspace{0.2cm}
  Recommandation actuelle: pour path-frozen, garder souvent
  \texttt{local\_pair\_mode="geodesic"}, car le goulot principal est plutôt
  l'optimisation répétée que Dijkstra lui-même.
\end{frame}

\begin{frame}{Path-frozen: accélération GPU}
  Le GPU n'exécute pas Dijkstra.

  \vspace{0.2cm}
  Il accélère surtout les évaluations répétées pendant L-BFGS-B:
  \begin{itemize}
    \item longueurs de toutes les arêtes des chemins gelés;
    \item accumulation des longueurs par chemin;
    \item résidus et gradients;
    \item paires locales directes si elles existent.
  \end{itemize}

  \vspace{0.2cm}
  Pourquoi Dijkstra reste CPU:
  \begin{itemize}
    \item structure très irrégulière, files de priorité, branchements;
    \item moins naturel pour le calcul vectorisé massif;
    \item pas le goulot dominant quand beaucoup d'itérations internes sont faites.
  \end{itemize}
\end{frame}

\begin{frame}{Soft-Bellman-Ford}
  Principe: remplacer le minimum de Bellman-Ford par un softmin.

  \vspace{0.2cm}
  Pour une relaxation:
  \[
    d_{t+1}(s,j)
    =
    \operatorname{softmin}_\beta
    \left(
      d_t(s,j),
      \{d_t(s,i)+c_{ij}\}_{i\to j}
    \right).
  \]

  \vspace{0.2cm}
  Avantage: les distances sont différentiables par rapport aux coûts d'arêtes,
  donc par rapport aux coordonnées.

  \vspace{0.2cm}
  Limite: contrairement à path-frozen, chaque évaluation de l'objectif doit
  refaire les relaxations, car les soft distances changent avec les coûts.
\end{frame}

\begin{frame}{Soft-BF: techniques pour le rendre utilisable}
  \begin{itemize}
    \item \textbf{sources par landmarks}: lancer BF depuis quelques sources seulement;
    \item \textbf{source\_batch\_size}: batcher les sources sur GPU pour contrôler la mémoire;
    \item \textbf{n\_relaxations}: tronquer les chemins à un nombre de sauts utile;
    \item \textbf{beta}: contrôler le compromis distance lisse / shortest-path dur;
    \item \textbf{local\_pair\_mode="direct"}: traiter les voisins locaux par distance directe;
    \item \textbf{on\_unreachable="warn\_skip"}: ignorer les paires non atteintes avec warning;
    \item \textbf{prob\_dtype=float32}: stocker la trace de backprop en précision réduite.
  \end{itemize}
\end{frame}

\begin{frame}{Soft-BF: rôle de \texorpdfstring{$\beta$}{beta}}
  La distance soft peut être vue comme
  \[
    d_\beta(A,B)
    =
    -\frac{1}{\beta}
    \log
    \sum_{\gamma:A\to B}
    e^{-\beta L(\gamma)} .
  \]

  \vspace{0.2cm}
  \begin{itemize}
    \item petit $\beta$: gradients diffus, objectif lisse mais biais entropique fort;
    \item grand $\beta$: gradients concentrés, meilleur shortest-path, objectif moins lisse;
    \item trop grand trop tôt: L-BFGS-B peut devenir lent ou instable.
  \end{itemize}

  \vspace{0.2cm}
  Heuristique naturelle: continuation en $\beta$, par exemple
  \[
    15 \rightarrow 25 \rightarrow 50 .
  \]
\end{frame}

\begin{frame}{DataSP / soft-Floyd-Warshall}
  Principe: soft Floyd-Warshall calcule directement des distances all-pairs
  différentiables.

  \vspace{0.2cm}
  Intérêt conceptuel:
  \begin{itemize}
    \item formulation très proche du softmin sur tous les chemins;
    \item pas de choix explicite d'un chemin gelé;
    \item utile comme référence pour petits graphes.
  \end{itemize}

  \vspace{0.2cm}
  Limite pratique:
  \begin{itemize}
    \item boucle cubique en $n$ dans la structure Floyd-Warshall;
    \item mémoire et trace de gradient très coûteuses;
    \item peu viable directement pour $n\simeq 3000$ sans autre approximation.
  \end{itemize}
\end{frame}

\begin{frame}{Masques de paires: éviter le all-pairs}
  Le all-pairs complet active $n(n-1)$ contraintes.

  \vspace{0.2cm}
  Stratégie utilisée:
  \begin{itemize}
    \item \textbf{landmarks}: contraintes globales depuis quelques sources;
    \item \textbf{local neighbors}: préserver la géométrie locale de $\Dtarget$;
    \item \textbf{random pairs}: ajouter un échantillon global non structuré;
    \item \textbf{max targets per source}: contrôler le coût sans perdre les sources.
  \end{itemize}

  \vspace{0.2cm}
  Cette stratégie sépare deux rôles:
  \begin{itemize}
    \item les landmarks imposent la structure globale;
    \item les voisins locaux empêchent les déchirures locales.
  \end{itemize}
\end{frame}

\begin{frame}{Local direct vs local géodésique}
  Pour une paire locale $(i,j)$:

  \vspace{0.2cm}
  \begin{block}{Mode géodésique}
    On impose que la distance le long du graphe reproduise $\Dtarget_{ij}$.
    C'est cohérent avec l'objectif géodésique, mais coûteux si tous les points
    deviennent sources.
  \end{block}

  \begin{block}{Mode direct}
    On impose directement $\Fmetric(x_j-x_i)\simeq \Dtarget_{ij}$.
    C'est une approximation tangentielle locale de la variété.
  \end{block}

  \vspace{0.1cm}
  Usage actuel:
  \begin{itemize}
    \item path-frozen: plutôt géodésique;
    \item soft-BF: plutôt direct, pour éviter BF depuis tous les sommets.
  \end{itemize}
\end{frame}

\begin{frame}{Gestion des paires inaccessibles}
  Si \texttt{n\_relaxations} est trop faible, ou si le graphe est mal connecté,
  certaines paires actives peuvent rester inaccessibles.

  \vspace{0.2cm}
  Deux comportements:
  \begin{itemize}
    \item \texttt{on\_unreachable="raise"}: mode strict, utile pour débugger;
    \item \texttt{on\_unreachable="warn\_skip"}: mode par défaut, ignore les paires concernées.
  \end{itemize}

  \vspace{0.2cm}
  Interprétation:
  \begin{itemize}
    \item quelques paires ignorées: acceptable pour un masque landmark large;
    \item beaucoup de paires ignorées: augmenter \texttt{n\_neighbors},
      \texttt{n\_relaxations} ou revoir le masque.
  \end{itemize}
\end{frame}

\begin{frame}{Pipeline pratique recommandé}
  \begin{enumerate}
    \item \textbf{Warm-start}: UMAP, Isomap ou Randers-SMACOF.
    \item \textbf{Path-frozen}: corriger rapidement la géométrie géodésique.
    \item \textbf{Soft-BF}: finisher plus lisse, avec $\beta$ élevé mais budget contrôlé.
  \end{enumerate}

  \vspace{0.3cm}
  Exemple de finisher soft-BF:
  \begin{itemize}
    \item \texttt{landmark\_mode="sources"};
    \item \texttt{n\_landmarks} assez grand mais pas all-pairs;
    \item \texttt{local\_pair\_mode="direct"};
    \item \texttt{target\_sampling="mixed"};
    \item \texttt{max\_targets\_per\_source} modéré;
    \item \texttt{beta} grand, éventuellement par continuation.
  \end{itemize}
\end{frame}

\begin{frame}{Messages à retenir}
  \begin{itemize}
    \item Randers-SMACOF est rapide et utile comme initialisation globale.
    \item Path-frozen transforme un problème discret en surrogate différentiable efficace.
    \item Soft-BF évite les chemins gelés mais paie les relaxations à chaque évaluation.
    \item Les masques landmarks + voisins locaux sont essentiels au passage à l'échelle.
    \item Le GPU est utile pour les évaluations massives, pas pour tous les sous-problèmes.
    \item Le bon finisher n'est pas forcément all-pairs: il doit équilibrer global et local.
  \end{itemize}
\end{frame}

\end{document}
